% GENERATED FILE. Edit canonical lessons, then run npm run generate:books.
% canonical-source-hash: fnv1a64:0e7822bd9ad5d00d
% canonical-lessons: ML-C505-urjjam, ML-C505-ksamam, ML-C505-bhugarbhajalam, ML-C505-agolatapanam, ML-C505-malakkuli

\chapter{Saving Resources}
\label{ch:ml-saving-resources}
\hlchaptermodality{505}

\begin{chapteropening}
\textbf{By the end of this chapter:} \emph{I can say energy, a shortage, groundwater, global warming and a rain pit.}

Complete the last lesson of chapter 505: I can say energy, a shortage, groundwater, global warming and a rain pit.
\end{chapteropening}

\section[\texorpdfstring{\ml{ഊർജ്ജം} (ūrjjaṁ)}{ഊർജ്ജം (ūrjjaṁ)}]{\ml{ഊർജ്ജം} (ūrjjaṁ) — energy}

\label{lesson:ML-C505-urjjam}



\begin{quote}
Before the new one: say the Malayalam for to be ruined, then the Malayalam for to reuse.
\end{quote}

\subsection*{You'll want to know: \ml{ഊർജ്ജം}}
\textbf{\ml{ഊർജ്ജം}} — \emph{ūrjjaṁ} — \textquotedblleft{}energy\textquotedblright{}.

\begin{quote}
\ml{സൂര്യനിൽ} \ml{നിന്ന്} \ml{നമുക്ക്} \ml{ഊർജ്ജം} \ml{കിട്ടും}.
\end{quote}

\begin{quote}
\emph{sūryanil ninnŭ namukkŭ ūrjjaṁ kiṭṭuṁ.}
\end{quote}

\begin{quote}
We can get energy from the sun.
\end{quote}

Written with a doubled \ml{ജ്ജ} after \ml{ർ}, the way \ml{ഓർമ്മ} doubles its \ml{മ}.

\begin{cousinweb}
\textbf{\ml{ഊർജ്ജം}} (\emph{ūrjjaṁ}) is borrowed from Sanskrit \emph{ūrja}, \textquotedblleft{}strong, powerful\textquotedblright{} (Monier-Williams).
\end{cousinweb}

\begin{grammarlens}[title={What you've built}]
One more word for getting around.
\end{grammarlens}

\subsection*{Your turn}
\begin{itemize}
\raggedright
  \item \emph{Say it:} \emph{ūrjjaṁ}
  \item \emph{Say it:} \emph{ūrjjaṁ}, once more
  \item \emph{Say it:} \emph{sūryanil ninnŭ namukkŭ ūrjjaṁ kiṭṭuṁ.} — We can get energy from the sun.
  \item \emph{Recall:} say the Malayalam for to be ruined, then the Malayalam for to reuse, then say \emph{ūrjjaṁ} again
\end{itemize}

\begin{tcolorbox}[breakable,colback=teal!4,colframe=teal!35!black,title={Before you move on}]
What does \ml{ഊർജ്ജം} mean? (\textquotedblleft{}Energy\textquotedblright{}.) Say it once more.
\end{tcolorbox}
\section[\texorpdfstring{\ml{ക്ഷാമം} (kṣāmaṁ)}{ക്ഷാമം (kṣāmaṁ)}]{\ml{ക്ഷാമം} (kṣāmaṁ) — a shortage; a scarcity}

\label{lesson:ML-C505-ksamam}



\begin{quote}
Before the new one: say the Malayalam for to reuse, then the Malayalam for energy.
\end{quote}

\subsection*{You'll want to know: \ml{ക്ഷാമം}}
\textbf{\ml{ക്ഷാമം}} — \emph{kṣāmaṁ} — \textquotedblleft{}a shortage; a scarcity\textquotedblright{}.

\begin{quote}
\ml{വേനലിൽ} \ml{ഗ്രാമത്തിൽ} \ml{കുടിവെള്ളത്തിന്} \ml{ക്ഷാമമുണ്ട്}.
\end{quote}

\begin{quote}
\emph{vēnalil grāmattil kuṭiveḷḷattinŭ kṣāmamuṇṭŭ.}
\end{quote}

\begin{quote}
In the dry season the village is short of drinking water.
\end{quote}

What is short takes -\ml{ന്} / -\ml{ക്ക്}: \ml{വെള്ളത്തിന്} \ml{ക്ഷാമം}, a shortage of water.

\begin{cousinweb}
\textbf{\ml{ക്ഷാമം}} (\emph{kṣāmaṁ}) is borrowed from Sanskrit \emph{kṣāma}, \textquotedblleft{}scorched, emaciated\textquotedblright{} (Monier-Williams).
\end{cousinweb}

\begin{grammarlens}[title={What you've built}]
One more word for getting around.
\end{grammarlens}

\subsection*{Your turn}
\begin{itemize}
\raggedright
  \item \emph{Say it:} \emph{kṣāmaṁ}
  \item \emph{Say it:} \emph{kṣāmaṁ}, once more
  \item \emph{Say it:} \emph{vēnalil grāmattil kuṭiveḷḷattinŭ kṣāmamuṇṭŭ.} — In the dry season the village is short of drinking water.
  \item \emph{Recall:} say the Malayalam for to reuse, then the Malayalam for energy, then say \emph{kṣāmaṁ} again
\end{itemize}

\begin{tcolorbox}[breakable,colback=teal!4,colframe=teal!35!black,title={Before you move on}]
What does \ml{ക്ഷാമം} mean? (\textquotedblleft{}A shortage; a scarcity\textquotedblright{}.) Say it once more.
\end{tcolorbox}
\section[\texorpdfstring{\ml{ഭൂഗർഭജലം} (bhūgarbhajalaṁ)}{ഭൂഗർഭജലം (bhūgarbhajalaṁ)}]{\ml{ഭൂഗർഭജലം} (bhūgarbhajalaṁ) — groundwater}

\label{lesson:ML-C505-bhugarbhajalam}



\begin{quote}
Before the new one: say the Malayalam for energy, then the Malayalam for a shortage.
\end{quote}

\subsection*{You'll want to know: \ml{ഭൂഗർഭജലം}}
\textbf{\ml{ഭൂഗർഭജലം}} — \emph{bhūgarbhajalaṁ} — \textquotedblleft{}groundwater\textquotedblright{}.

\begin{quote}
\ml{വേനലിൽ} \ml{ഭൂഗർഭജലം} \ml{താഴുന്നു}.
\end{quote}

\begin{quote}
\emph{vēnalil bhūgarbhajalaṁ tāḻunnu.}
\end{quote}

\begin{quote}
In the dry season the groundwater level falls.
\end{quote}

\ml{ഭൂ}-, the earth, plus \ml{ഗർഭം}, the womb, plus \ml{ജലം}, water: water in the belly of the earth.

\begin{cousinweb}
\textbf{\ml{ഭൂഗർഭജലം}} (\emph{bhūgarbhajalaṁ}) is a Malayalam compound built from three Sanskrit words, \emph{bhū}, \textquotedblleft{}the earth\textquotedblright{}, \emph{garbha}, \textquotedblleft{}the womb, the inside\textquotedblright{}, and \emph{jala}, \textquotedblleft{}water\textquotedblright{} (Monier-Williams).
\end{cousinweb}

\begin{grammarlens}[title={What you've built}]
One more word for getting around.
\end{grammarlens}

\subsection*{Your turn}
\begin{itemize}
\raggedright
  \item \emph{Say it:} \emph{bhūgarbhajalaṁ}
  \item \emph{Say it:} \emph{bhūgarbhajalaṁ}, once more
  \item \emph{Say it:} \emph{vēnalil bhūgarbhajalaṁ tāḻunnu.} — In the dry season the groundwater level falls.
  \item \emph{Recall:} say the Malayalam for energy, then the Malayalam for a shortage, then say \emph{bhūgarbhajalaṁ} again
\end{itemize}

\begin{tcolorbox}[breakable,colback=teal!4,colframe=teal!35!black,title={Before you move on}]
What does \ml{ഭൂഗർഭജലം} mean? (\textquotedblleft{}Groundwater\textquotedblright{}.) Say it once more.
\end{tcolorbox}
\section[\texorpdfstring{\ml{ആഗോളതാപനം} (āgōḷatāpanaṁ)}{ആഗോളതാപനം (āgōḷatāpanaṁ)}]{\ml{ആഗോളതാപനം} (āgōḷatāpanaṁ) — global warming}

\label{lesson:ML-C505-agolatapanam}



\begin{quote}
Before the new one: say the Malayalam for a shortage, then the Malayalam for groundwater.
\end{quote}

\subsection*{You'll want to know: \ml{ആഗോളതാപനം}}
\textbf{\ml{ആഗോളതാപനം}} — \emph{āgōḷatāpanaṁ} — \textquotedblleft{}global warming\textquotedblright{}.

\begin{quote}
\ml{ആഗോളതാപനം} \ml{കാരണം} \ml{കടലിലെ} \ml{വെള്ളം} \ml{ഉയരുന്നു}.
\end{quote}

\begin{quote}
\emph{āgōḷatāpanaṁ kāraṇaṁ kaṭalile veḷḷaṁ uyarunnu.}
\end{quote}

\begin{quote}
Because of global warming the sea is rising.
\end{quote}

\ml{ആഗോള}, global, from \ml{ഗോളം}, a globe, plus \ml{താപനം}, heating.

\begin{cousinweb}
\textbf{\ml{ആഗോളതാപനം}} (\emph{āgōḷatāpanaṁ}) is a modern Malayalam compound built from Sanskrit \emph{gola}, \textquotedblleft{}a ball, a globe\textquotedblright{}, and \emph{tāpana}, \textquotedblleft{}burning\textquotedblright{} (Monier-Williams), which Malayalam uses for heating.
\end{cousinweb}

\begin{grammarlens}[title={What you've built}]
One more word for getting around.
\end{grammarlens}

\subsection*{Your turn}
\begin{itemize}
\raggedright
  \item \emph{Say it:} \emph{āgōḷatāpanaṁ}
  \item \emph{Say it:} \emph{āgōḷatāpanaṁ}, once more
  \item \emph{Say it:} \emph{āgōḷatāpanaṁ kāraṇaṁ kaṭalile veḷḷaṁ uyarunnu.} — Because of global warming the sea is rising.
  \item \emph{Recall:} say the Malayalam for a shortage, then the Malayalam for groundwater, then say \emph{āgōḷatāpanaṁ} again
\end{itemize}

\begin{tcolorbox}[breakable,colback=teal!4,colframe=teal!35!black,title={Before you move on}]
What does \ml{ആഗോളതാപനം} mean? (\textquotedblleft{}Global warming\textquotedblright{}.) Say it once more.
\end{tcolorbox}
\section[\texorpdfstring{\ml{മഴക്കുഴി} (maḻakkuḻi)}{മഴക്കുഴി (maḻakkuḻi)}]{\ml{മഴക്കുഴി} (maḻakkuḻi) — a rain pit}

\label{lesson:ML-C505-malakkuli}



\begin{quote}
Before the new one: say the Malayalam for groundwater, then the Malayalam for global warming.
\end{quote}

\subsection*{You'll want to know: \ml{മഴക്കുഴി}}
\textbf{\ml{മഴക്കുഴി}} — \emph{maḻakkuḻi} — \textquotedblleft{}a rain pit\textquotedblright{}.

\begin{quote}
\ml{വീടിന്റെ} \ml{പിന്നിൽ} \ml{ഞങ്ങൾ} \ml{ഒരു} \ml{മഴക്കുഴി} \ml{കുഴിച്ചു}.
\end{quote}

\begin{quote}
\emph{vīṭinṟe pinnil ñaṅṅaḷ oru maḻakkuḻi kuḻiccu.}
\end{quote}

\begin{quote}
We dug a rain pit behind the house.
\end{quote}

A pit dug to catch rainwater so that it sinks into the ground and fills the wells.

\begin{cousinweb}
\textbf{\ml{മഴക്കുഴി}} (\emph{maḻakkuḻi}) is a native Dravidian compound of \ml{മഴ}, rain (DEDR 4753), and \ml{കുഴി}, a pit (DEDR 1818).
\end{cousinweb}

\begin{grammarlens}[title={What you've built}]
One more word for getting around.
\end{grammarlens}

\subsection*{Your turn}
\begin{itemize}
\raggedright
  \item \emph{Say it:} \emph{maḻakkuḻi}
  \item \emph{Say it:} \emph{maḻakkuḻi}, once more
  \item \emph{Say it:} \emph{vīṭinṟe pinnil ñaṅṅaḷ oru maḻakkuḻi kuḻiccu.} — We dug a rain pit behind the house.
  \item \emph{Recall:} say the Malayalam for groundwater, then the Malayalam for global warming, then say \emph{maḻakkuḻi} again
\end{itemize}

\begin{tcolorbox}[breakable,colback=teal!4,colframe=teal!35!black,title={Before you move on}]
What does \ml{മഴക്കുഴി} mean? (\textquotedblleft{}A rain pit\textquotedblright{}.) Say it once more.
\end{tcolorbox}
